\subsection{Runtime Robustness and Source Reachability}\label{sec:method:runtime}

To study runtime robustness, {\tech} executes bytecode-level inputs under multiple CPython versions, including CPython 3.10, 3.11, and 3.12, and optionally debug or sanitizer builds.  Inputs include real bytecode artifacts, structured bytecode variants, and a raw byte-mutation baseline.  Outcomes are classified as load failure, Python-level exception, normal termination, timeout, assertion failure, abort, segmentation fault, or sanitizer finding.  We deduplicate abnormal outcomes using signal type, top stack frame, failure location, assertion message, crash-address bucket, input hash, and code-object feature signatures.

For each abnormal outcome, {\tech} checks source reachability.  It attempts to decompile the bytecode, recompile the recovered source, re-execute under the same runtime, and compare the failure signature.  Outcomes are classified as source-reproducible, bytecode-only under current tools, decompilation failed, recompilation failed, behavior mismatch, or inconclusive.  This step separates bytecode-level robustness findings from ordinary source-level vulnerabilities and identifies where source-based testing cannot reproduce the behavior observed at the bytecode boundary.

This stage is included because robustness results without reachability are easy to overstate.  Raw bytecode mutation can generate artifacts that no normal compiler would emit, while ordinary source-level testing may miss behaviors that valid loaders can still execute.  Reporting source reachability alongside crash class keeps these interpretations separate.
